\section{Implications (Epistemic Only)}
\label{sec:implications}

In this paper, the term ``implications'' refers exclusively to \emph{epistemic
consequences} of the formal setting.
These consequences constrain what can be inferred, asserted, or ruled out given
a fixed observation interface and a fixed vocabulary boundary
(\texttt{ETS.K\_EA.v1.1}).
They do not constitute operational, managerial, normative, or methodological
recommendations.

\subsection{Implications for Inference Under Incomplete Observability}
\label{subsec:imp-incomplete-observability}

The formal definitions and results entail the following epistemic limitations:

\begin{enumerate}
  \item \textbf{Non-uniqueness is the default.}
  In the absence of explicit identifiability conditions, the inverse relation
  from observations to internal structure is set-valued.
  Multiple mutually incompatible internal realizations may remain
  observationally indistinguishable.

  \item \textbf{Apparent stability may be epistemic rather than structural.}
  Stable patterns in observations do not entail stability of the underlying
  internal structure.
  Such apparent stability may result from observational coarseness, masking
  effects, or information loss at the observation interface.

  \item \textbf{Additional data does not guarantee improved diagnosability.}
  Supplementary observations may be epistemically irrelevant if they do not
  refine the indistinguishability classes induced by the observation operator.
  Claims of improved diagnosability therefore require formal justification
  relative to the model and interface, not intuitive appeal.

  \item \textbf{Boundary effects dominate near critical regions.}
  In the vicinity of phase boundaries, including STOP, small changes in
  observations or assumptions may induce discontinuous changes in admissible
  inferences.
  This sensitivity reflects epistemic instability of inference, not necessarily
  volatility of the underlying system.
\end{enumerate}

\subsection{Implications for the Meaning of Phase Regions}
\label{subsec:imp-phase-regions}

A phase-region label derived from the predicates defined in this paper must be
interpreted as follows:

\begin{itemize}
  \item \textbf{A classification within the formal semantics.}
  Phase regions are defined relative to explicit primitives, assumptions, and
  the declared observation interface.
  They have no meaning outside this formal context.

  \item \textbf{Not an assessment of desirability.}
  A region is neither favorable nor unfavorable in a managerial, operational,
  or normative sense.
  It expresses only diagnosability and existence conditions within the model.

  \item \textbf{Not a guarantee of recoverability or control.}
  Membership in a non-STOP region does not imply that the system is recoverable,
  controllable, stabilizable, or improvable under any intervention scheme.
\end{itemize}

\subsection{Implications of STOP}
\label{subsec:imp-stop}

The STOP classification has the following epistemic implications:

\begin{enumerate}
  \item \textbf{Inference collapse.}
  Under the given observation interface and formal assumptions, non-degenerate
  structural diagnosis is not epistemically justified.

  \item \textbf{Constraint violation or undecidability.}
  STOP indicates either violation of existence or consistency constraints, or a
  loss of diagnosability required to support meaningful structural claims within
  the admitted vocabulary.

  \item \textbf{Non-substitutability.}
  STOP cannot be replaced by any positive diagnostic classification without
  introducing additional assumptions, new primitives, or semantic extensions
  beyond the declared formal boundary.
\end{enumerate}

\subsection{Implications for Claims and Reporting}
\label{subsec:imp-claims-reporting}

The formal results impose negative constraints on what can be claimed in a
scientifically valid report:

\begin{enumerate}
  \item \textbf{No causal claims without identifiability.}
  When multiple incompatible internal realizations are compatible with the same
  observations, causal assertions about internal structure are epistemically
  underdetermined.

  \item \textbf{No operational conclusions from epistemic classifications.}
  Translating a diagnostic classification into a decision, recommendation, or
  intervention requires an explicit policy, control, or optimization model,
  none of which is provided or assumed here.

  \item \textbf{No universality without a declared mapping.}
  Any application to a concrete enterprise must explicitly specify the mapping
  from empirical artifacts and observations to the formal objects of the model.
  Without such a mapping, reported results lack determinate interpretive content.
\end{enumerate}
